\chapter{Gallbladder}

\label{ch:ch46_gallbladder}

Ultrasound is the best first test for gallbladder problems. It is quick, painless, uses no radiation, and can detect gallstones with very high accuracy. During an ultrasound of the gallbladder, the sonographer will press the probe firmly over the right side of your abdomen, just below the ribs. This may cause some discomfort, especially if your gallbladder is inflamed. The sonographer will ask you to take deep breaths to help visualise the gallbladder.

Gallstones are extremely common---about 10--15\% of adults have them, and most cause no symptoms. When a gallstone blocks the cystic duct, it causes \textbf{acute cholecystitis} (inflammation of the gallbladder), which presents with right upper quadrant pain, fever, and tenderness. If a stone migrates into the common bile duct (\textbf{choledocholithiasis}), it can cause jaundice and pancreatitis and requires more advanced imaging such as MRCP (a special MRI of the bile ducts).

Your doctor may also request an ultrasound to evaluate a gallbladder polyp (a growth on the gallbladder wall) or to assess the bile ducts for dilation.


\section{First-Line Assessment}
Ultrasound commonly evaluates gallstones, cholecystitis, bile-duct dilation, and gallbladder wall abnormalities. CT and MRI with MRCP provide complementary assessment when ultrasound is limited or complications are suspected.

\begin{table}[h]
\centering
\caption{Gallbladder imaging modalities}
\begin{tabular}{p{3cm} p{4cm} p{4cm}}
\hline
\textbf{Modality} & \textbf{Strengths} & \textbf{Limitations} \\
\hline
CT & Evaluates alternative diagnoses and complications; useful when ultrasound is equivocal or the clinical question requires cross-sectional imaging & May miss noncalcified gallstones; not usually the initial test for uncomplicated biliary colic \\
ERCP & Therapeutic (stone extraction, stenting) & Invasive; risk of pancreatitis; reserved for intervention \\
EUS & Highest sensitivity for small CBD stones ($< 5$ mm) & Semi-invasive; operator-dependent \\
MRCP & Non-invasive evaluation of bile ducts; excellent for choledocholithiasis & Expensive; limited availability; takes 30--45 min \\
Ultrasound & Usually the initial study for suspected gallstones or acute cholecystitis; high accuracy for gallbladder stones & Operator-dependent; limited by body habitus and bowel gas; distal CBD may be obscured \\
\hline
\end{tabular}
\end{table}

\begin{pearlbox}[Ultrasound for biliary pain]
For suspected gallstones or acute cholecystitis presenting with right upper quadrant pain, ultrasound is usually the initial imaging study. Further imaging, including cholescintigraphy, CT, or MRI/MRCP, depends on the suspected diagnosis, ultrasound result, and need to assess complications or alternatives.
\end{pearlbox}


\section{Ultrasound Review of the Gallbladder}

\subsection{Technique}
The patient should be fasting for at least 6 hours to allow the gallbladder to fill with bile (fasting distends the gallbladder, making stones easier to detect). The probe is placed in the right upper quadrant in longitudinal and subcostal oblique planes.

\subsection{Key Measurements and Findings}
\begin{itemize}
  \item \textbf{Gallbladder wall thickness}: Normal $< 3$ mm when distended. Wall thickening $> 3$ mm suggests cholecystitis, but can also occur in fasting, hepatitis, heart failure, and hypoalbuminaemia.
  \item \textbf{Gallbladder size}: Normal length 7--10 cm, transverse diameter $< 4$ cm. Distension $> 4$ cm suggests obstruction.
  \item \textbf{Pericholecystic fluid}: Fluid around the gallbladder wall---a sign of acute cholecystitis.
  \item \textbf{Stones}: Echogenic foci with posterior acoustic shadowing; they layer dependently. Small stones may not shadow.
  \item \textbf{Sludge}: Low-level echoes without shadowing; gravity-dependent.
  \item \textbf{Common bile duct (CBD) diameter}: Normal $< 6$ mm (or $< 8$ mm post-cholecystectomy, and increases 1 mm per decade after age 60). CBD dilation suggests choledocholithiasis.
\end{itemize}

\begin{pearlbox}[Sonographic Murphy sign]
The \textbf{sonographic Murphy sign} is focal tenderness elicited by probe pressure directly over the visualised gallbladder. It supports acute cholecystitis when combined with other imaging, clinical, and laboratory findings, but it is not definitive and may be absent, including after analgesia.
\end{pearlbox}


\section{Acute Cholecystitis}
\textbf{Ultrasound diagnostic criteria for acute cholecystitis}:
\begin{enumerate}
  \item Gallstones (present in 90--95\% of cases)
  \item Sonographic Murphy sign (positive in 80--90\%)
  \item Gallbladder wall thickening ($> 3$ mm) (present in 50--70\%)
  \item Pericholecystic fluid (present in 20--50\%)
  \item Gallbladder distension ($> 4$ cm)
\end{enumerate}

Gallstones with a sonographic Murphy sign and inflammatory findings increase suspicion for acute cholecystitis but do not establish the diagnosis in isolation. If ultrasound is equivocal and the result will affect management, cholescintigraphy (HIDA) may help assess cystic duct obstruction; interpret it with clinical context.

\subsection{Complications of Acute Cholecystitis}
\begin{table}[h]
\centering
\caption{Complications of acute cholecystitis}
\begin{tabular}{p{3cm} p{4.5cm} p{4cm}}
\hline
\textbf{Complication} & \textbf{Imaging features} & \textbf{Significance} \\
\hline
Emphysematous cholecystitis & Gas in gallbladder wall or lumen on CT & More common in diabetics; requires urgent surgery \\
Empyema & Pus in gallbladder lumen; thick enhancing wall; marked pericholecystic inflammation & Urgent decompression (percutaneous cholecystostomy or surgery) \\
Gangrenous cholecystitis & Intraluminal membranes (sloughed mucosa), asymmetric wall thickening, pericholecystic abscess & Requires emergency cholecystectomy \\
Perforation & Pericholecystic fluid collection, discontinuity of gallbladder wall, free fluid & Emergent \\
\hline
\end{tabular}
\end{table}


\section{Choledocholithiasis}
Choledocholithiasis (stones in the common bile duct) presents with jaundice, raised bilirubin, and dilated CBD.

\begin{itemize}
  \item \textbf{Ultrasound}: Dilated CBD $> 6$ mm; echogenic focus with shadowing in the CBD (sensitivity 50--75\%---CBD stones are frequently missed on ultrasound because the distal CBD is obscured by duodenal gas).
  \item \textbf{MRCP}: Noninvasive assessment for ductal stones and obstruction; performance varies with stone size and protocol, and very small stones may be missed.
  \item \textbf{EUS}: Highest sensitivity for small stones ($< 5$ mm) not seen on MRCP.
  \item \textbf{ERCP}: Both diagnostic and therapeutic---stone extraction with sphincterotomy.
\end{itemize}

\begin{pearlbox}[When to order MRCP]
Order MRCP when:
\begin{itemize}
  \item Ultrasound shows dilated CBD or echogenic CBD but stone not confirmed
  \item Liver function tests are cholestatic (raised ALP, GGT, conjugated bilirubin)
  \item Clinical suspicion of choledocholithiasis despite negative ultrasound
  \item Pre-ERCP planning (to confirm stone burden and plan intervention)
\end{itemize}
Do NOT order MRCP as a screening test---it should be targeted to a specific clinical question.
\end{pearlbox}


\section{Cholangitis}
Ascending cholangitis is infection of the biliary tree, typically due to obstruction (most commonly choledocholithiasis).

\textbf{Clinical triad (Charcot's triad)}: Fever, jaundice, right upper quadrant pain.

\textbf{Imaging}:
\begin{itemize}
  \item \textbf{Ultrasound}: Dilated CBD, may show stone; sometimes periductal edema
  \item \textbf{CT}: Dilated ducts, wall enhancement, periductal fluid; may show the obstructing cause
  \item \textbf{MRCP}: Best non-invasive assessment of the biliary tree; shows obstruction level and cause
  \item \textbf{ERCP}: Therapeutic---drainage and stone extraction
\end{itemize}

\begin{pearlbox}[Reynold's pentad]
Charcot's triad (fever, jaundice, RUQ pain) + hypotension and altered mental status = \textbf{Reynold's pentad}, indicating suppurative (life-threatening) cholangitis requiring emergency biliary drainage (ERCP or percutaneous transhepatic cholangiography).
\end{pearlbox}


\section{Gallbladder Polyps}
Incidental gallbladder polyps are common. Management depends on size:

\begin{table}[h]
\centering
\caption{Gallbladder polyp management}
\begin{tabular}{p{3cm} p{4cm} p{4cm}}
\hline
\textbf{Size} & \textbf{Risk of malignancy} & \textbf{Management} \\
\hline
6--9 mm & Low (1--3\%) & Ultrasound follow-up at 6 months, then annually for 2 years \\
$\geq 10$ mm & Moderate to high (up to 30--40\%) & Cholecystectomy recommended \\
$< 6$ mm & Very low ($< 1\%$) & Ultrasound follow-up at 12 months; if stable, no further follow-up \\
 sessile polyp, rapid growth, or symptomatic & High & Cholecystectomy regardless of size \\
\hline
\end{tabular}
\end{table}

\begin{pearlbox}[Adenomyomatosis vs polyp]
Adenomyomatosis is a benign condition with Rokitansky--Aschoff sinuses (intramural diverticula) that appear as echogenic foci with ``comet-tail'' artefact on ultrasound. Unlike polyps, adenomyomatosis does NOT require cholecystectomy unless symptomatic.
\end{pearlbox}


\section{Gallbladder Cancer}
Gallbladder carcinoma is rare but aggressive. Imaging features:
\begin{itemize}
  \item \textbf{Ultrasound}: Focal or diffuse gallbladder wall thickening; mass replacing the gallbladder; intraluminal mass. May be incidentally found after cholecystectomy for ``cholelithiasis.''
  \item \textbf{CT}: Mass replacing the gallbladder; direct invasion into liver (segments IV and V); lymphadenopathy; peritoneal deposits.
  \item \textbf{Risk factors}: Large gallstones ($> 3$ cm), chronic cholecystitis, porcelain gallbladder, gallbladder polyps $\geq 10$ mm.
  \item \textbf{Staging}: TNM staging guides surgical planning (radical cholecystectomy vs palliative care).
\end{itemize}


\section{Biliary Dilation}
Dilated intrahepatic ducts appear as ``parallel channel sign'' (dilated bile duct running alongside the portal vein, giving a ``double barrel shotgun'' appearance on ultrasound).

\textbf{Causes of biliary dilation}:
\begin{itemize}
  \item \textbf{Obstructive}: Choledocholithiasis (most common), pancreatic head mass, cholangiocarcinoma (Klatskin tumor at hilum), ampullary tumor
  \item \textbf{Non-obstructive}: Post-surgical strictures, primary sclerosing cholangitis, Caroli disease
\end{itemize}

The level of obstruction is critical:
\begin{itemize}
  \item \textbf{Intrahepatic only}: Think primary sclerosing cholangitis, infiltration, or high strictures
  \item \textbf{Extrahepatic}: Think choledocholithiasis, pancreatic mass, or distal cholangiocarcinoma
  \item \textbf{Hilar (Klatskin tumor)}: Bismuth classification guides surgical approach
\end{itemize}


\section{Review}
Assess stones, sludge, wall thickness, distension, sonographic tenderness, pericholecystic fluid, hyperemia, bile ducts, liver, and adjacent bowel. Interpret wall findings with fasting state, systemic illness, and technique.

\section{Escalation}
Report suspected gangrene, perforation, emphysematous infection, obstructing stone, ascending infection, or biliary malignancy promptly. Do not use a single equivocal sign as a substitute for clinical assessment.

\section*{Selected Guidance}
\begin{itemize}
    \item American College of Radiology. \href{https://acsearch.acr.org/docs/69474/Narrative/}{Appropriateness Criteria: Right Upper Quadrant Pain}.
    \item World Society of Emergency Surgery. \href{https://pmc.ncbi.nlm.nih.gov/articles/PMC7643471/}{Guidelines for diagnosis and treatment of acute calculus cholecystitis}.
\end{itemize}

\clearpage
